% =====================================================================
% VERSION ENSAE, REDUCTION SFE (phase 3, 1er octobre 2026) : supplement S3, questionnaire d'elicitation reproduit (ancien C.3).
% Copie a texte constant de chapitres/18_elicitation_protocole.tex, lignes 97 a 252.
% Ne changent que les titres et les renvois repares, listes dans
% reduction/phase3.md.
% Contenu RETIRE de la version ENSAE et destine au supplement en ligne : aucun
% fichier maitre ne le lit. Voir supplement/README.md.
% =====================================================================

% Declaration reprise telle quelle de chapitres/18_elicitation_protocole.tex (ligne 40),
% qui couvrait ce questionnaire quand il vivait dans l'annexe :
% HARNAIS-HORS-SCRIPT: les quantiles demandes a l'expert sont une propriete de CONCEPTION du
% questionnaire, non une grandeur calculee : aucun script ne les imprime et aucun ne le devrait.

%--------------------------------------------------------------
\chapter{Questionnaire d'élicitation}
\label{sec:questionnaire-reproduit}
%--------------------------------------------------------------
% LES VALEURS VRAIES DES GRAINES NE SONT PAS PUBLIEES ICI, ET C'EST DELIBERE : les imprimer
% BRULERAIT l'instrument, un expert qui a lu le memoire ne pouvant plus etre note dessus. Le
% script 105 les calcule, leur applique quatre controles et les imprime dans sa sortie de
% travail ; si l'elicitation etait lancee apres publication, les graines se retirent du meme jeu
% de donnees par ce meme script. Ce qui est cite ci-dessous est donc le seul nombre que
% l'arbitrage de conception exige de montrer, la valeur du candidat ecarte.
% SOURCES-SCRIPTS: 105

Décrire un instrument ne suffit pas à établir qu'il existe. Il est donc reproduit ici dans la
forme sous laquelle il aurait été soumis, de façon qu'un lecteur puisse juger sur pièce ce que
la section précédente affirme, et le réutiliser si la décision était rouverte. Le formulaire
complet, avec sa consigne, son exemple traité, sa mention d'usage et de conservation et son
accord signé du répondant, est versionné à part dans le dépôt du mémoire ; la présente annexe en
reprend les deux tables et l'essentiel de la consigne.

\begin{attention}
\textbf{L'instrument reproduit ici est la version révisée, et il diffère de l'exemplaire de
travail.} Celui-ci portait une question d'étalonnage sur l'indice de queue de sévérité, retirée
pour les deux motifs exposés plus haut. Elle n'est pas corrigée mais \emph{remplacée}, et le
remplacement a lui-même été arbitré : un indice de Gini de la concentration avait d'abord été
retenu, puis écarté parce qu'il fait doublon avec la question qui compte les incidents portant la
moitié du volume, et parce qu'une valeur vraie de $0{,}97$ se devine sans rien connaître du
dossier. La graine retenue porte sur la durée de confinement, qui est une performance de
résilience et donc du même domaine que les cibles. L'exemplaire de travail ne doit pas
être soumis en l'état.
\end{attention}

\paragraph{La règle que s'impose l'instrument, vérifiée graine par graine.} Une graine doit être une
\emph{grandeur réalisée}, reproductible par une sortie de script versionnée, et jamais un
paramètre estimé. Les dix graines retenues sont donc des comptages, des parts, des délais ou des
quantiles observés : aucune ne demande d'ajuster un modèle. Un script les calcule et leur applique
quatre contrôles, portant sur ce caractère réalisé, sur la non-dégénérescence de la valeur, sur le
respect des bornes naturelles, et sur le pouvoir de séparation entre experts. Une seule graine est
signalée comme potentiellement devinable, celle qui porte sur la part du piratage ; elle est
conservée avec son motif : sa valeur vraie est nettement au-dessus de la réponse spontanée, si
bien qu'elle sépare malgré son caractère extrême.

\paragraph{Aucune graine ne porte une grandeur de capital, choix délibéré.} Un praticien de la
résilience observe des incidents, des délais et des concentrations ; il n'observe pas un quantile
de charge annuelle. Noter un expert sur une telle grandeur mesurerait sa familiarité avec le
présent mémoire, non son jugement sur le risque.

\section*{Consigne remise à l'expert}

\begin{cle}
\textbf{Comment répondre.} Pour chaque quantité, donnez trois nombres : une valeur \emph{basse}
($5\,\%$), une valeur \emph{médiane} ($50\,\%$), une valeur \emph{haute} ($95\,\%$), telles que
vous estimez à $90\,\%$ de chances que la vraie valeur soit comprise entre votre $5\,\%$ et votre
$95\,\%$. Il n'y a pas de réponse attendue : donnez votre meilleure estimation \emph{et} votre
incertitude. Répondez même incertain, un intervalle large étant une réponse honnête et
parfaitement valable. Ne cherchez pas les vraies valeurs : l'objet est de mesurer votre jugement.
\end{cle}

\section*{Partie A, les dix questions d'étalonnage}

Des quantités observées dans les bases du mémoire, dont la valeur vraie est connue de l'analyste
et n'est pas affichée au répondant. Elles servent à pondérer les experts, non à les piéger. Les
trois colonnes de droite sont vides dans les deux tables qui suivent : c'est l'espace de réponse
du formulaire, non une donnée manquante.

{\small
\renewcommand{\arraystretch}{1.22}
\begin{longtable}{@{}l p{8.0cm} c c c@{}}
\toprule
& \textbf{Quantité, sur les incidents notifiés de la chronologie} & $\mathbf{5\,\%}$ &
$\mathbf{50\,\%}$ & $\mathbf{95\,\%}$ \\
\midrule
\endfirsthead
\toprule
& \textbf{Quantité (suite)} & $\mathbf{5\,\%}$ & $\mathbf{50\,\%}$ & $\mathbf{95\,\%}$ \\
\midrule
\endhead
G1 & Délai médian entre la survenance d'un incident et sa déclaration, en jours & & & \\
G2 & Part des incidents déclarés plus de quatre-vingt-dix jours après leur survenance,
     en \% & & & \\
G3 & Délai médian de déclaration d'un incident de type piratage, en jours & & & \\
G4 & Durée médiane entre le début et la fin d'exposition d'un incident, en jours & & & \\
G5 & Nombre médian d'enregistrements touchés par incident & & & \\
G6 & Nombre d'incidents qui portent à eux seuls la moitié du volume touché
     sur sept exercices & & & \\
G7 & Part des incidents dont l'exposition a duré plus de trente jours, en \% & & & \\
G8 & Part des incidents de type piratage parmi ceux dont le vecteur est renseigné,
     en \% & & & \\
G9 & Part des lignes de la chronologie qui sont des doublons de notification d'un même
     incident, en \% & & & \\
G10 & Rapport du quatre-vingt-dix-neuvième centile à la médiane des pertes cyber du
      secteur financier & & & \\
\bottomrule
\end{longtable}}

\vspace{2pt}
\noindent\textbf{Ce que chaque graine cherche à mesurer.} G1 à G3 portent sur le \emph{délai de
détection}, qui est le canal que la conformité déplace le plus directement. G4 et G7 portent sur
le \emph{confinement}. G5, G6 et G10 portent sur la \emph{concentration} de la perte, c'est-à-dire
sur l'intuition de queue lourde. G8 porte sur la structure des vecteurs et G9 sur la qualité de la
donnée déclarative, dont un praticien de la notification a une expérience directe.

\section*{Partie B, les six liens dirigés entre piliers}

Force du lien dirigé de $j$ vers $k$, entre $0$ et $1$ : $0$ signifie qu'une défaillance du
pilier $j$ n'entraîne jamais celle de $k$, $1$ qu'elle l'entraîne quasi certainement dans la
foulée. Les trois mêmes quantiles sont demandés.

\begin{attention}
\textbf{Le sens des lettres n'est pas une coquette\-rie de notation.} La matrice du modèle
s'écrit $W_{jk}$ avec $j$ la \emph{source} et $k$ la \emph{cible} : la ligne émet, la colonne
reçoit, et deux scripts vérifient cette convention avant tout calcul. Un formulaire qui
demanderait la force du lien \emph{de $k$ vers $j$} ferait donc saisir la matrice
\emph{transposée}. Or c'est précisément entre une matrice et sa transposée que le
chapitre~\ref{ch:identifiabilite} montre que la donnée ne tranche pas, et le
chapitre~\ref{chap:cascade} que le capital et la décision diffèrent. Une inversion de lettres
dans l'instrument produirait donc exactement l'erreur que tout le mémoire s'emploie à ne pas
commettre, et elle ne serait rattrapée par aucun contrôle numérique.
\end{attention}

{\small
\renewcommand{\arraystretch}{1.25}
\begin{longtable}{@{}l p{8.2cm} c c c@{}}
\toprule
& \textbf{Lien dirigé $j \rightarrow k$, $j$ émet} & $\mathbf{5\,\%}$ & $\mathbf{50\,\%}$ &
$\mathbf{95\,\%}$ \\
\midrule
\endfirsthead
\toprule
& \textbf{Lien dirigé (suite)} & $\mathbf{5\,\%}$ & $\mathbf{50\,\%}$ & $\mathbf{95\,\%}$ \\
\midrule
\endhead
B1 & P1 $\rightarrow$ P2 : une défaillance de gouvernance entraîne des incidents & & & \\
B2 & P2 $\rightarrow$ P1 : des incidents entraînent une défaillance de gouvernance & & & \\
B3 & P4 $\rightarrow$ P2 : une défaillance d'un tiers entraîne des incidents & & & \\
B4 & P1 $\rightarrow$ P4 : une défaillance de gouvernance entraîne une défaillance côté
     tiers & & & \\
B5 & P3 $\rightarrow$ P1 : une défaillance des tests de résilience entraîne une défaillance
     de gouvernance & & & \\
B6 & P5 $\rightarrow$ P2 : une défaillance du partage d'information entraîne des
     incidents & & & \\
\bottomrule
\end{longtable}}

\paragraph{Le couple B1 et B2 est l'objet même du mémoire.} Les deux questions portent sur le
\emph{même} couple de piliers, dans les deux sens. Une donnée de co-occurrence ne les distingue
pas ; un expert, lui, peut les séparer. C'est précisément la quantité que le
chapitre~\ref{ch:identifiabilite} montre non identifiable : ce
questionnaire aurait donc eu une valeur que nulle autre source du dossier ne porte.

\paragraph{Ce qui est fait des réponses.} La partie~A mesure la calibration du répondant, c'est-à-dire
la justesse de ses intervalles, et son information, c'est-à-dire leur finesse ; ces deux
grandeurs fixent son poids dans l'agrégation. La partie~B, pondérée par ce poids et croisée avec
les réponses des autres experts, produirait chaque lien de la matrice sous forme d'\emph{intervalle}
et non de point. Le traitement est automatique et aucune donnée d'incident n'est requise, seulement
le jugement.
